%% cap01-planteamiento-unesco.tex — Planteamiento del problema (diseño UNESCO, exploratorio)
%% Sustituye a cap01-planteamiento.tex (eje confuciano) cuando se integre a main.tex.

\providecommand{\pendiente}[1]{\textbf{[Pendiente: #1]}}

\chapter{Planteamiento del problema}
\label{cap:planteamiento}

\section{Situación de la práctica profesional}
\label{sec:practica}

Mi trabajo consiste en diseñar currículo de cursos sobre inteligencia artificial en
DeepLearning.AI, una organización de educación en línea sobre esta tecnología. En paralelo
estoy por concluir la edición 2026 del Programa regional de formación ``Planeamiento y
gestión de políticas educativas en América Latina y el Caribe'', del IIPE UNESCO
\citep{iipe2026prf}, cuyos cursos cubren el análisis situacional, la formulación, la
implementación y el monitoreo de políticas educativas.

Las dos actividades se tocan en un punto. Quien diseña un currículo de inteligencia
artificial decide qué deben aprender las personas sobre esta tecnología, y esas decisiones
no ocurren en el vacío: las políticas nacionales ya declaran qué competencias esperan, para
quién y con qué propósito. Quien analiza esas políticas, a su vez, necesita saber qué dicen,
qué callan y en qué se parecen a las de otros países. Mi práctica de diseño curricular me
coloca como usuario de esas políticas; mi formación en el IIPE, como su analista.

En esa práctica encuentro una dificultad concreta. El análisis de políticas que aprendo y
aplico descansa en la lectura de documentos completos, uno por uno, con categorías que el
analista define para cada ejercicio. Funciona bien para un documento o para dos, pero no
escala. Cuando la pregunta es cómo tratan siete países un mismo tema, en cuatro idiomas y
con documentos de decenas o cientos de páginas, la lectura manual obliga a elegir: o se
reduce el número de casos, o se reduce la profundidad, o se lee todo en traducción al
inglés. Cada una de esas salidas cambia el resultado antes de empezar.

La situación que quiero transformar es esa forma de analizar políticas públicas: pasar de
una lectura manual, limitada a pocos casos y a una lengua, a un procedimiento que pueda
aplicarse con criterios fijos a muchos documentos en su idioma original, y cuyos resultados
otra persona pueda revisar y repetir.

\section{Situación problemática}
\label{sec:problematica}

La inteligencia artificial llegó a la agenda educativa de los gobiernos en pocos años.
UNESCO convocó en 2019 una conferencia internacional sobre el tema cuyo documento final, el
Consenso de Beijing, recomienda a los Estados planear la inteligencia artificial en sus
políticas educativas \citep{unesco2019beijing}, y en 2021 publicó una guía para que los
responsables de política diseñaran esas políticas \citep{miao2021guidance}. Desde entonces
los países han publicado estrategias, planes de acción e informes que comparten buena parte
de su vocabulario.

Ese vocabulario común plantea el problema de esta tesis. Cuando dos países escriben sobre
``competencias para la era de la inteligencia artificial'' o sobre ``uso ético de los
datos'', no es evidente si dicen lo mismo, si repiten la fórmula de un organismo
internacional o si usan las mismas palabras para propósitos distintos. Los estudios que
han comparado estas políticas lo muestran: la convergencia en el vocabulario convive con
una divergencia en la interpretación y en la puesta en práctica
\citep{jobin2019landscape, fjeld2020principled}, y cuando las estrategias hablan de
educación, suelen referirse a formar fuerza laboral más que a transformar la enseñanza
\citep{schiff2022education}.

La forma en que se han hecho esas comparaciones limita lo que se puede concluir de ellas.
El estado del arte (sección~\ref{subsec:ea-sintesis}) muestra cuatro rasgos recurrentes. Se
comparan pocos casos, entre dos y cuatro países. Se trabaja con documentos en inglés o
traducidos al inglés, y los propios autores reconocen el sesgo que eso introduce. Las
categorías de comparación se derivan del corpus o las construyen los autores, de modo que
cada estudio usa una rejilla distinta. Y la codificación la hace una sola persona, sin
medida de acuerdo, como reconoce \citet{schiff2022education}.

México vive esta situación desde afuera. La primera propuesta de estrategia nacional de
inteligencia artificial fue una agenda de 2018 impulsada fuera del gobierno, que no se
convirtió en política pública \citep{cminds2018iamexico}. En abril de 2026 la SEP presentó
líneas de acción para el uso ético y crítico de la inteligencia artificial generativa en la
educación superior, con carácter de recomendación \citep{sep2026ia}, y la ANUIES publicó en
2024 un informe sobre gobernanza de la inteligencia artificial \citep{anuies2024gobernanza}.
Pero el país no tiene todavía un documento de política nacional que pueda compararse con los
de otros países; por eso México queda fuera del corpus, y esa ausencia se reporta como
resultado. Cuando ese documento exista, hará falta un procedimiento para compararlo con lo
que otros países ya publicaron.

\section{Preguntas de investigación}
\label{sec:preguntas}

\subsection{Pregunta central}

Cuando las políticas de educación en inteligencia artificial de siete países se leen con las
siete categorías de la guía de UNESCO para responsables de política, ¿qué tan cercanas o
lejanas están entre sí dentro de cada categoría, y qué muestra esa distancia sobre la forma
en que cada país traduce un vocabulario internacional común?

\subsection{Preguntas derivadas}

\begin{enumerate}
  \item[\textbf{P1.}] \textbf{Cobertura.} ¿Qué categorías de la guía de UNESCO atiende cada
  política y cuáles omite, y cuánto peso tiene en cada una la educación \textit{para} la
  inteligencia artificial frente a la inteligencia artificial \textit{para} la educación?

  \item[\textbf{P2.}] \textbf{Distancia.} Dentro de una misma categoría, ¿qué países quedan
  cerca y cuáles lejos en un espacio semántico multilingüe, y ese acomodo cambia de una
  categoría a otra?

  \item[\textbf{P3.}] \textbf{Robustez.} ¿Las distancias entre países se sostienen cuando se
  controla el idioma de los documentos y cuando se considera el grado de acuerdo del panel de
  modelos que los clasifica?
\end{enumerate}

P1 describe qué dice cada política, P2 compara cómo lo dice y P3 pone a prueba si la
comparación depende del idioma o del instrumento.

\section{Objetivos}
\label{sec:objetivos}

\subsection{Objetivo general}

Comparar, de manera exploratoria y reproducible, las políticas de educación en inteligencia
artificial de siete países dentro de cada una de las siete categorías de la guía de UNESCO,
mediante un procedimiento automático de procesamiento de lenguaje natural y un panel de
modelos de lenguaje multilingües.

\subsection{Objetivos específicos}

\begin{enumerate}
  \item[\textbf{OE1.}] Construir un corpus de documentos oficiales de política con criterios
  de inclusión explícitos, conservados en su idioma original y divididos en fragmentos
  recuperables por país.
  \item[\textbf{OE2.}] Clasificar cada fragmento, con su contexto inmediato, en las siete
  categorías de UNESCO y en la distinción entre educación para la inteligencia artificial e
  inteligencia artificial para la educación, mediante un panel de siete modelos de lenguaje
  de distinto origen (responde a P1).
  \item[\textbf{OE3.}] Estimar, dentro de cada categoría, la distancia semántica entre los
  fragmentos de cada par de países (responde a P2).
  \item[\textbf{OE4.}] Contrastar las distancias con un control de idioma y con el acuerdo
  entre los modelos del panel (responde a P3).
  \item[\textbf{OE5.}] Proponer una herramienta abierta que permita a otros analistas aplicar
  el mismo procedimiento a nuevos documentos de política.
\end{enumerate}

\section{Supuestos}
\label{sec:supuestos}

El estudio es exploratorio, por lo que no plantea hipótesis que se acepten o rechacen. Parte
de tres supuestos que se declaran para que puedan discutirse.

\begin{description}
  \item[S1.] Los siete apartados de la sección 6 de la guía de UNESCO pueden usarse como
  categorías de comparación, aunque la guía no los presente así
  (sección~\ref{subsec:mt-marcos}).
  \item[S2.] La cercanía entre fragmentos en un espacio semántico multilingüe refleja, al
  menos en parte, cercanía de contenido y no solo de idioma. P3 pone a prueba este supuesto.
  \item[S3.] Un panel de modelos de distinto origen, con categorías fijadas de antemano,
  reduce el riesgo de que la clasificación dependa de la elección de un solo modelo
  \citep{baumann2025llmhacking}.
\end{description}

\section{Justificación}
\label{sec:justificacion}

La justificación pedagógica está en el objeto. Las políticas de educación en inteligencia
artificial deciden qué se enseña sobre la tecnología, cómo se usa en las aulas y qué se
espera de los docentes. Quien diseña currículo trabaja dentro de esas decisiones, y
compararlas con las de otros países permite ver que no son las únicas posibles.

La justificación metodológica está en el procedimiento. Los estudios revisados comparan
pocos casos, en inglés y con rejillas propias (sección~\ref{subsec:ea-sintesis}). Esta tesis
compara siete países, en su idioma original y con una rejilla tomada de un organismo
internacional, y en lugar de un solo codificador usa un panel de modelos cuyo acuerdo se
mide. El procedimiento sigue la postura de \citet{grimmer2013text}: los métodos automáticos
amplían la capacidad del analista pero no sustituyen la interpretación, y sus resultados
deben validarse.

La justificación práctica está en México. Cuando el país publique su política de
inteligencia artificial en educación, contar con un procedimiento abierto para compararla
con la de otros países permitirá discutirla con evidencia y no solo con referencias
generales a lo que recomiendan los organismos internacionales.

\section{Alcances y limitaciones}
\label{sec:alcances}

El corpus reúne 15 documentos de siete países: Alemania, Australia, Canadá, China,
Colombia, Estados Unidos y Sudáfrica, publicados entre 2020 y 2026, en alemán, chino,
español e inglés, divididos en 3{,}095 fragmentos. China aporta nueve documentos; el resto de
los países, uno cada uno. El 15.º Plan Quinquenal chino, por su extensión, se excluye de las
comparaciones de distancia para que China no pese más que los demás países, y esa exclusión
se declara en cada resultado.

El estudio es exploratorio. Describe cómo se parecen y en qué se distinguen las políticas; no
prueba por qué. Analiza lo que los documentos dicen, no si los gobiernos lo cumplen: opera en
los planos del texto y del discurso, no en el de los efectos \citep{rizvi2010globalizing}.
Con siete casos no hay generalización estadística posible. En cuatro de los países
(Alemania, Australia, Canadá y Estados Unidos) la educación depende en buena medida de los
estados o provincias, de modo que el documento nacional no equivale a la política que rige
en sus escuelas \citep{bray1995levels}.

Hay una limitación que la rejilla misma introduce. Si los países redactaron sus políticas
siguiendo las recomendaciones de UNESCO, medirlas con las categorías de UNESCO podría
producir convergencia por construcción. Por eso la tesis no mide cuánto se apega cada país a
la guía, sino cuán lejos quedan los países entre sí dentro de cada categoría: dos países
pueden atender la misma categoría y hacerlo de maneras muy distintas.

La tesis se inspira en la investigación-acción, pero no incluye una intervención. Su aporte
a la práctica es un procedimiento y una herramienta abierta para analizar políticas; su uso
con otros analistas ocurrirá después de concluido este estudio.
